% Build: xelatex pfx.slides.tex
%
% Converted from pfx.slides.odt (LibreOffice Writer, not Impress --
% five pages separated by manual page breaks standing in for slides)
% to a real beamer deck. xelatex, for the Hebrew in the author line.
\documentclass[11pt]{beamer}
\setbeamertemplate{navigation symbols}{}
\usepackage{pfxmacros}
\usepackage{fontspec}
\newfontfamily\hebrewfont{David CLM}[Script=Hebrew]
\TeXXeTstate=1

\title{Prefix Macros for THE and XEDIT}
\author{Seymour J. Metz (H: Shmuel \texorpdfstring{{\hebrewfont\beginR שמואל בן לייביש ולאה\endR}}{שמואל בן לייביש ולאה})}
\institute{\href{mailto:smetz3@gmu.edu}{smetz3@gmu.edu} \\ \url{http://mason.gmu.edu/~smetz3}}
\date{}

\begin{document}

\begin{frame}
\titlepage
\end{frame}

\begin{frame}{Agenda}
\begin{itemize}
\item Background
\item What it does
\item Demo of revised version
\item How it works
\end{itemize}
\end{frame}

\begin{frame}{Background}
\begin{itemize}
\item Large document dissemination system
\item Front end ran under CMS in VM/SP HPO
\item Written in
  \begin{itemize}
  \item Assembler H with PROC macros
  \item PL/I
  \item Several scripting languages
  \end{itemize}
\item Documentation in DCF GML SS
\item Home grown IDE
\item No syntax oriented editor
\item The prefix macros were an extension to the IDE
\end{itemize}
\end{frame}

\begin{frame}{What it does}
\begin{itemize}
\item Inserts templates for control structures, with empty lines for additional input
\item Automatically indents
\item Sets cursor appropriately for indentation structure
\item Includes some utility commands, e.g., UC
\item Commands support repetition counts and other parameters, although seriously limited by the size of the prefix field
\end{itemize}
\end{frame}

\begin{frame}{Demo of revised version}
\end{frame}

\begin{frame}{How it works}
\begin{itemize}
\item Requires a \constant{SET PREFIX SYNONYM} for each command
  \begin{itemize}
  \item Typically done in profile
  \end{itemize}
\item Makes heavy use of \constant{SET PENDING}
\item Defines and make heavy use of
  \begin{itemize}
  \item \constant{..}
  \item \constant{RI}
  \item \constant{RIC}
  \end{itemize}
\item Common housekeeping to, e.g., extract an optional count parameter
\end{itemize}
\end{frame}

\end{document}
